\documentclass[12pt,a4paper]{article}
\usepackage[T2A]{fontenc}
\usepackage[utf8]{inputenc}
\usepackage[russian]{babel}
\usepackage{amsmath,amssymb,geometry,hyperref}
\geometry{margin=18mm}
\setlength{\parindent}{0pt}
\setlength{\parskip}{5pt}
\begin{document}
\begin{center}\large Тема 1. Вариант 4\end{center}
Борисюк Е. Г.
\par\textbf{Условие}\[
z=3x_1+x_2+4x_3+2x_4\to\min,\qquad
\begin{cases}
2x_1+x_2+x_3\le8,\\
x_1+x_3+x_4=5,\\
x_2+x_4\ge4,\\
x_i\ge0.
\end{cases}
\]
\par\textbf{Канонический вид}
Добавляю запас $s_1$ в первое ограничение, вычитаю избыток $s_3$ из третьего. Второе уже является равенством.
\[
2x_1+x_2+x_3+s_1=8,\quad
x_1+x_3+x_4=5,\quad
x_2+x_4-s_3=4.
\]
\par\textbf{Вспомогательная задача}
Для начального базиса добавляю $a_2,a_3$.
\[
w=a_2+a_3\to\min,\qquad
\begin{cases}
2x_1+x_2+x_3+s_1=8,\\
x_1+x_3+x_4+a_2=5,\\
x_2+x_4-s_3+a_3=4.
\end{cases}
\]
Все переменные неотрицательны. Начальный базис:
$s_1=8,\ a_2=5,\ a_3=4$, остальные переменные равны нулю.

Оценки: $r_j=c_j-\sum_i c_{B_i}t_{ij}$.
Выбираю первый столбец с $r_j<0$, строку по минимуму $b_i/t_{ik}$ при $t_{ik}>0$.
\par\textbf{Таблица I$_{0}$}\begin{center}\small\setlength{\tabcolsep}{4pt}\begin{tabular}{c|rrrrrrrr|r}\hline
Базис & $x_{1}$ & $x_{2}$ & $x_{3}$ & $x_{4}$ & $s_{1}$ & $s_{3}$ & $a_{2}$ & $a_{3}$ & $b$\\\hline
$s_{1}$ & $2$ & $1$ & $1$ & $0$ & $1$ & $0$ & $0$ & $0$ & $8$\\
$a_{2}$ & $1$ & $0$ & $1$ & $1$ & $0$ & $0$ & $1$ & $0$ & $5$\\
$a_{3}$ & $0$ & $1$ & $0$ & $1$ & $0$ & $-1$ & $0$ & $1$ & $4$\\
\hline $r_j$ & $-1$ & $-1$ & $-1$ & $-2$ & $0$ & $1$ & $0$ & $0$ & $9$\\\hline\end{tabular}\end{center}
\[\theta_{1}=8/(2)=4,\quad \theta_{2}=5/(1)=5.\]
$x_{1}$ входит, $s_{1}$ выходит. Разрешающий элемент $2$.
\[R_{1}'=R_{1}/(2),\quad R_{2}'=R_{2}-(1)R_{1}'.\]
\newpage
\par\textbf{Таблица I$_{1}$}\begin{center}\small\setlength{\tabcolsep}{4pt}\begin{tabular}{c|rrrrrrrr|r}\hline
Базис & $x_{1}$ & $x_{2}$ & $x_{3}$ & $x_{4}$ & $s_{1}$ & $s_{3}$ & $a_{2}$ & $a_{3}$ & $b$\\\hline
$x_{1}$ & $1$ & $\frac{1}{2}$ & $\frac{1}{2}$ & $0$ & $\frac{1}{2}$ & $0$ & $0$ & $0$ & $4$\\
$a_{2}$ & $0$ & $\frac{-1}{2}$ & $\frac{1}{2}$ & $1$ & $\frac{-1}{2}$ & $0$ & $1$ & $0$ & $1$\\
$a_{3}$ & $0$ & $1$ & $0$ & $1$ & $0$ & $-1$ & $0$ & $1$ & $4$\\
\hline $r_j$ & $0$ & $\frac{-1}{2}$ & $\frac{-1}{2}$ & $-2$ & $\frac{1}{2}$ & $1$ & $0$ & $0$ & $5$\\\hline\end{tabular}\end{center}
\[\theta_{1}=4/(\frac{1}{2})=8,\quad \theta_{3}=4/(1)=4.\]
$x_{2}$ входит, $a_{3}$ выходит. Разрешающий элемент $1$.
\[R_{3}'=R_{3}/(1),\quad R_{1}'=R_{1}-(\frac{1}{2})R_{3}',\quad R_{2}'=R_{2}-(\frac{-1}{2})R_{3}'.\]
\par\textbf{Таблица I$_{2}$}\begin{center}\small\setlength{\tabcolsep}{4pt}\begin{tabular}{c|rrrrrrrr|r}\hline
Базис & $x_{1}$ & $x_{2}$ & $x_{3}$ & $x_{4}$ & $s_{1}$ & $s_{3}$ & $a_{2}$ & $a_{3}$ & $b$\\\hline
$x_{1}$ & $1$ & $0$ & $\frac{1}{2}$ & $\frac{-1}{2}$ & $\frac{1}{2}$ & $\frac{1}{2}$ & $0$ & $\frac{-1}{2}$ & $2$\\
$a_{2}$ & $0$ & $0$ & $\frac{1}{2}$ & $\frac{3}{2}$ & $\frac{-1}{2}$ & $\frac{-1}{2}$ & $1$ & $\frac{1}{2}$ & $3$\\
$x_{2}$ & $0$ & $1$ & $0$ & $1$ & $0$ & $-1$ & $0$ & $1$ & $4$\\
\hline $r_j$ & $0$ & $0$ & $\frac{-1}{2}$ & $\frac{-3}{2}$ & $\frac{1}{2}$ & $\frac{1}{2}$ & $0$ & $\frac{1}{2}$ & $3$\\\hline\end{tabular}\end{center}
\[\theta_{1}=2/(\frac{1}{2})=4,\quad \theta_{2}=3/(\frac{1}{2})=6.\]
$x_{3}$ входит, $x_{1}$ выходит. Разрешающий элемент $\frac{1}{2}$.
\[R_{1}'=R_{1}/(\frac{1}{2}),\quad R_{2}'=R_{2}-(\frac{1}{2})R_{1}'.\]
\newpage
\par\textbf{Таблица I$_{3}$}\begin{center}\small\setlength{\tabcolsep}{4pt}\begin{tabular}{c|rrrrrrrr|r}\hline
Базис & $x_{1}$ & $x_{2}$ & $x_{3}$ & $x_{4}$ & $s_{1}$ & $s_{3}$ & $a_{2}$ & $a_{3}$ & $b$\\\hline
$x_{3}$ & $2$ & $0$ & $1$ & $-1$ & $1$ & $1$ & $0$ & $-1$ & $4$\\
$a_{2}$ & $-1$ & $0$ & $0$ & $2$ & $-1$ & $-1$ & $1$ & $1$ & $1$\\
$x_{2}$ & $0$ & $1$ & $0$ & $1$ & $0$ & $-1$ & $0$ & $1$ & $4$\\
\hline $r_j$ & $1$ & $0$ & $0$ & $-2$ & $1$ & $1$ & $0$ & $0$ & $1$\\\hline\end{tabular}\end{center}
\[\theta_{2}=1/(2)=\frac{1}{2},\quad \theta_{3}=4/(1)=4.\]
$x_{4}$ входит, $a_{2}$ выходит. Разрешающий элемент $2$.
\[R_{2}'=R_{2}/(2),\quad R_{1}'=R_{1}-(-1)R_{2}',\quad R_{3}'=R_{3}-(1)R_{2}'.\]
\par\textbf{Таблица I$_{4}$}\begin{center}\small\setlength{\tabcolsep}{4pt}\begin{tabular}{c|rrrrrrrr|r}\hline
Базис & $x_{1}$ & $x_{2}$ & $x_{3}$ & $x_{4}$ & $s_{1}$ & $s_{3}$ & $a_{2}$ & $a_{3}$ & $b$\\\hline
$x_{3}$ & $\frac{3}{2}$ & $0$ & $1$ & $0$ & $\frac{1}{2}$ & $\frac{1}{2}$ & $\frac{1}{2}$ & $\frac{-1}{2}$ & $\frac{9}{2}$\\
$x_{4}$ & $\frac{-1}{2}$ & $0$ & $0$ & $1$ & $\frac{-1}{2}$ & $\frac{-1}{2}$ & $\frac{1}{2}$ & $\frac{1}{2}$ & $\frac{1}{2}$\\
$x_{2}$ & $\frac{1}{2}$ & $1$ & $0$ & $0$ & $\frac{1}{2}$ & $\frac{-1}{2}$ & $\frac{-1}{2}$ & $\frac{1}{2}$ & $\frac{7}{2}$\\
\hline $r_j$ & $0$ & $0$ & $0$ & $0$ & $0$ & $0$ & $1$ & $1$ & $0$\\\hline\end{tabular}\end{center}
$w^*=0$, искусственные переменные равны нулю и выведены из базиса.
\newpage
\par\textbf{Основная задача}

$w^*=0$, поэтому исходная задача допустима. Удаляю столбцы $a_2,a_3$.
Получен план $(0,7/2,9/2,1/2)$.
Возвращаю $c=(3,1,4,2,0,0)$; $c_B=(4,2,1)$.
Пересчитываю оценки; $z=45/2$.
\par\textbf{Таблица II$_{0}$}\begin{center}\small\setlength{\tabcolsep}{4pt}\begin{tabular}{c|rrrrrr|r}\hline
Базис & $x_{1}$ & $x_{2}$ & $x_{3}$ & $x_{4}$ & $s_{1}$ & $s_{3}$ & $b$\\\hline
$x_{3}$ & $\frac{3}{2}$ & $0$ & $1$ & $0$ & $\frac{1}{2}$ & $\frac{1}{2}$ & $\frac{9}{2}$\\
$x_{4}$ & $\frac{-1}{2}$ & $0$ & $0$ & $1$ & $\frac{-1}{2}$ & $\frac{-1}{2}$ & $\frac{1}{2}$\\
$x_{2}$ & $\frac{1}{2}$ & $1$ & $0$ & $0$ & $\frac{1}{2}$ & $\frac{-1}{2}$ & $\frac{7}{2}$\\
\hline $r_j$ & $\frac{-5}{2}$ & $0$ & $0$ & $0$ & $\frac{-3}{2}$ & $\frac{-1}{2}$ & $\frac{45}{2}$\\\hline\end{tabular}\end{center}
\[\theta_{1}=\frac{9}{2}/(\frac{3}{2})=3,\quad \theta_{3}=\frac{7}{2}/(\frac{1}{2})=7.\]
$x_{1}$ входит, $x_{3}$ выходит. Разрешающий элемент $\frac{3}{2}$.
\[R_{1}'=R_{1}/(\frac{3}{2}),\quad R_{2}'=R_{2}-(\frac{-1}{2})R_{1}',\quad R_{3}'=R_{3}-(\frac{1}{2})R_{1}'.\]
\par\textbf{Таблица II$_{1}$}\begin{center}\small\setlength{\tabcolsep}{4pt}\begin{tabular}{c|rrrrrr|r}\hline
Базис & $x_{1}$ & $x_{2}$ & $x_{3}$ & $x_{4}$ & $s_{1}$ & $s_{3}$ & $b$\\\hline
$x_{1}$ & $1$ & $0$ & $\frac{2}{3}$ & $0$ & $\frac{1}{3}$ & $\frac{1}{3}$ & $3$\\
$x_{4}$ & $0$ & $0$ & $\frac{1}{3}$ & $1$ & $\frac{-1}{3}$ & $\frac{-1}{3}$ & $2$\\
$x_{2}$ & $0$ & $1$ & $\frac{-1}{3}$ & $0$ & $\frac{1}{3}$ & $\frac{-2}{3}$ & $2$\\
\hline $r_j$ & $0$ & $0$ & $\frac{5}{3}$ & $0$ & $\frac{-2}{3}$ & $\frac{1}{3}$ & $15$\\\hline\end{tabular}\end{center}
\[\theta_{1}=3/(\frac{1}{3})=9,\quad \theta_{3}=2/(\frac{1}{3})=6.\]
$s_{1}$ входит, $x_{2}$ выходит. Разрешающий элемент $\frac{1}{3}$.
\[R_{3}'=R_{3}/(\frac{1}{3}),\quad R_{1}'=R_{1}-(\frac{1}{3})R_{3}',\quad R_{2}'=R_{2}-(\frac{-1}{3})R_{3}'.\]
\newpage
\par\textbf{Таблица II$_{2}$}\begin{center}\small\setlength{\tabcolsep}{4pt}\begin{tabular}{c|rrrrrr|r}\hline
Базис & $x_{1}$ & $x_{2}$ & $x_{3}$ & $x_{4}$ & $s_{1}$ & $s_{3}$ & $b$\\\hline
$x_{1}$ & $1$ & $-1$ & $1$ & $0$ & $0$ & $1$ & $1$\\
$x_{4}$ & $0$ & $1$ & $0$ & $1$ & $0$ & $-1$ & $4$\\
$s_{1}$ & $0$ & $3$ & $-1$ & $0$ & $1$ & $-2$ & $6$\\
\hline $r_j$ & $0$ & $2$ & $1$ & $0$ & $0$ & $-1$ & $11$\\\hline\end{tabular}\end{center}
\[\theta_{1}=1/(1)=1.\]
$s_{3}$ входит, $x_{1}$ выходит. Разрешающий элемент $1$.
\[R_{1}'=R_{1}/(1),\quad R_{2}'=R_{2}-(-1)R_{1}',\quad R_{3}'=R_{3}-(-2)R_{1}'.\]
\par\textbf{Таблица II$_{3}$}\begin{center}\small\setlength{\tabcolsep}{4pt}\begin{tabular}{c|rrrrrr|r}\hline
Базис & $x_{1}$ & $x_{2}$ & $x_{3}$ & $x_{4}$ & $s_{1}$ & $s_{3}$ & $b$\\\hline
$s_{3}$ & $1$ & $-1$ & $1$ & $0$ & $0$ & $1$ & $1$\\
$x_{4}$ & $1$ & $0$ & $1$ & $1$ & $0$ & $0$ & $5$\\
$s_{1}$ & $2$ & $1$ & $1$ & $0$ & $1$ & $0$ & $8$\\
\hline $r_j$ & $1$ & $1$ & $2$ & $0$ & $0$ & $0$ & $10$\\\hline\end{tabular}\end{center}

Все оценки неотрицательны. Базис: $s_3=1,\ x_4=5,\ s_1=8$.
Небазисные $x_1=x_2=x_3=0$.
\[
\boxed{x^*=(0,0,0,5),\qquad z_{\min}=10.}
\]
Проверка: $0\le8$, $0+0+5=5$, $0+5\ge4$.
Excel Solver даёт тот же ответ.
\par\textbf{Рефлективный вывод}

Я разобрал построение начального базиса для равенства и ограничения
со знаком $\ge$. Наиболее трудным был пересчёт строк и оценок.
Для проверки сравнил таблицы с Python, а итоговый ответ с Excel Solver.
\end{document}
